\documentclass{szb-solution}

\area{Sorok}
\title{Numerikus sorok}
\subject{Matematika G1}
\subjectCode{BMETE94BG01}
\date{Utoljára frissítve: \today}
\docno{14}

\begin{document}
\maketitle

\subsection{Konvergencia a részletösszegek alapján}

\begin{enumerate}[a)]
  \item A sort két mértani sor összegére bontjuk:
        $$
          \sum_{n=1}^{\infty}\frac{2^n+3^n}{6^n}
          =\sum_{n=1}^{\infty}\left(\frac13\right)^n
           +\sum_{n=1}^{\infty}\left(\frac12\right)^n
          \text.
        $$
        A részletösszeg határértéke
        $$
          \frac{1/3}{1-1/3}+\frac{1/2}{1-1/2}
          =\boxed{\frac32}
          \text,
        $$
        tehát a sor konvergens.

  \item A feladatlapon a sor $n=1$-től indul, de az első tag nevezője ekkor
        nulla, tehát a sor így nincs értelmezve. A természetes javítás az
        $n=2$ kezdőindex. Ekkor
        $$
          \frac1{n(n-1)}=\frac1{n-1}-\frac1n
          \text,
        $$
        ezért a teleszkopikus részletösszeg
        $$
          s_N=\sum_{n=2}^{N}\left(\frac1{n-1}-\frac1n\right)
             =1-\frac1N\longrightarrow1
          \text.
        $$
        A javított sor összege \boxed{1}.

  \item A logaritmus azonosságai alapján
        \begin{align*}
          s_N
          &=\sum_{n=2}^{N}\ln\left(1-\frac1{n^2}\right)\\
          &=\ln\prod_{n=2}^{N}\frac{(n-1)(n+1)}{n^2}
            =\ln\frac{N+1}{2N}
          \text.
        \end{align*}
        Így
        $$
          \boxed{\sum_{n=2}^{\infty}\ln\left(1-\frac1{n^2}\right)
          =\ln\frac12=-\ln2}
          \text.
        $$

  \item Tegyük fel, hogy $a,k\in\mathbb R$. A
        $\sum a^n/n^k$ sor konvergenciája:
        \begin{itemize}
          \item $|a|<1$ esetén minden $k$ mellett abszolút konvergens;
          \item $|a|>1$ esetén divergens, mert a tagok nem tartanak nullához;
          \item $a=1$ esetén pontosan $k>1$ mellett konvergens;
          \item $a=-1$ esetén $k>1$ mellett abszolút, $0<k\leq1$ mellett
                feltételesen konvergens, $k\leq0$ mellett divergens.
        \end{itemize}
        Az $a=0$ esetet az első pont tartalmazza.
\end{enumerate}

\subsection{A Cauchy-féle konvergenciakritérium}

Legyen
$$
  a_n=\frac{n}{n^3+n^2+1}
  \text.
$$
Minden $n\geq1$ esetén $0<a_n\leq1/n^2$. Ha $q\geq p\geq2$, akkor
$$
  0\leq\sum_{n=p}^{q}a_n
  \leq\sum_{n=p}^{\infty}\frac1{n^2}
  \leq\int_{p-1}^{\infty}\frac{\dd x}{x^2}
  =\frac1{p-1}
  \text.
$$
Tetszőleges $\varepsilon>0$ esetén válasszuk
$p>1+1/\varepsilon$. Ekkor minden $q\geq p$ mellett a sorfarok kisebb
$\varepsilon$-nál, tehát a Cauchy-kritérium szerint a sor
\boxed{\text{konvergens}}.

\subsection{Konvergenciavizsgálat}

\begin{enumerate}[a)]
  \item Határértékes összehasonlítással
        $$
          \lim_{n\to\infty}
          \frac{(2n^3-16)/(n^5+n)}{1/n^2}=2
          \text.
        $$
        Mivel $\sum1/n^2$ konvergens, az adott sor is
        \boxed{\text{konvergens}}.

  \item $\cos(\pi/2)=0$, így minden tag nulla. A sor
        \boxed{\text{konvergens}}, összege $0$.

  \item $n\geq3$ esetén $\ln n<n$, ezért
        $1/\ln n>1/n$. A harmonikus sor divergenciája miatt az adott sor
        \boxed{\text{divergens}}.

  \item Ismert nevezetes határérték, hogy
        $$
          \left(1-\frac1n\right)^n\longrightarrow e^{-1}\neq0
          \text.
        $$
        A szükséges feltétel nem teljesül, tehát a sor
        \boxed{\text{divergens}}.

  \item Határértékes összehasonlítással
        $$
          \lim_{n\to\infty}
          \frac{1/\sqrt{n(n+1)}}{1/n}=1
          \text.
        $$
        Ezért a harmonikus sorral együtt \boxed{\text{divergens}}.

  \item A gyökteszt szerint
        $$
          \lim_{n\to\infty}
          \sqrt[n]{\frac{2n^2}{(2+1/n)^n}}
          =\frac12<1
          \text,
        $$
        tehát a sor \boxed{\text{konvergens}}.

  \item Az $n=1$ tag $0^0$ alakú, ezért azt külön kell megadni vagy a sort
        $n=2$-től indítani; egyetlen véges tag a konvergenciát nem befolyásolja.
        A gyökteszt határértéke
        $$
          \lim_{n\to\infty}
          \left(\frac{n-1}{n+1}\right)^{n-1}=e^{-2}<1
          \text,
        $$
        ezért a sor \boxed{\text{konvergens}}.

  \item A hányadosteszt alapján
        $$
          \lim_{n\to\infty}
          \frac{(n+1)/e^{n+1}}{n/e^n}=\frac1e<1
          \text.
        $$
        A sor \boxed{\text{konvergens}}; összege
        $e/(e-1)^2$.

  \item A $b_n=n/(n^2+1)$ sorozat nullához tart és $n\geq1$-től
        nemnövekvő. A Leibniz-kritérium miatt az alternáló sor konvergens.
        Ugyanakkor $b_n\sim1/n$, ezért abszolút értékben divergens. A sor
        \boxed{\text{feltételesen konvergens}}.

  \item Mivel
        $$
          \left|\frac{\sin n}{\sqrt[3]{n^4}}\right|
          \leq\frac1{n^{4/3}}
          \text,
        $$
        és a jobb oldali $p$-sor konvergens, az adott sor
        \boxed{\text{abszolút konvergens}}.

  \item Legyen $b_n=(n-1)/(n(n+1))$. Ekkor $b_n\to0$, és $n\geq2$-től
        nemnövekvő, ezért a Leibniz-kritérium konvergenciát ad. Viszont
        $$
          \lim_{n\to\infty}\frac{b_n}{1/n}=1
          \text,
        $$
        tehát az abszolútérték-sor divergens. A sor
        \boxed{\text{feltételesen konvergens}}.

  \item A feladatlapon ez a sor szó szerint megegyezik az előzővel, ezért az
        eredmény is ugyanaz: \boxed{\text{feltételesen konvergens}}.
\end{enumerate}

\subsection{Két sor összegsora}

Mindkét megadott sor abszolút konvergens, és tagonként
$$
  a_n+b_n
  =\frac{1+n+(-1)^n-n}{3^n}
  =\frac{1+(-1)^n}{3^n}
  \text.
$$
Ezért az összegsor konvergens, összege
$$
  \sum_{n=1}^{\infty}\frac1{3^n}
  +\sum_{n=1}^{\infty}\left(-\frac13\right)^n
  =\frac12-\frac14
  =\boxed{\frac14}
  \text.
$$

\subsection{Két divergens sor különbségsora}

Az $\sum1/(2n-1)$ és $\sum1/(2n)$ sor külön-külön divergens, de a
tagok különbsége
$$
  a_n-b_n
  =\frac1{2n-1}-\frac1{2n}
  =\frac1{2n(2n-1)}
  \text.
$$
Ez pozitív és legfeljebb $1/[2n(n-1)]$ ($n\geq2$), tehát a különbségsor
konvergens. Ez az alternáló harmonikus sor, így
$$
  \boxed{\sum_{n=1}^{\infty}(a_n-b_n)=\ln2}
  \text.
$$

\subsection{Cauchy-szorzat}

Az első sor
$$
  \sum_{n=1}^{\infty}\frac1{n\sqrt n}
  =\sum_{n=1}^{\infty}\frac1{n^{3/2}}
$$
abszolút konvergens $p$-sor, a második pedig abszolút konvergens mértani sor:
$$
  \sum_{n=1}^{\infty}\frac1{2^{n-1}}=2
  \text.
$$
Két abszolút konvergens sor Cauchy-szorzata abszolút konvergens, és összege
az összegek szorzata. Következésképpen
$$
  \boxed{\text{a Cauchy-szorzat konvergens, összege }
  2\zeta\!\left(\frac32\right)}
  \text.
$$

\end{document}
